\chapter{Теория}

Кластеризация — это набор методов обучения без учителя для группировки данных по определённым критериям в так называемые кластеры, что позволяет выявлять сходства и различия между объектами, а также упрощать их анализ и визуализацию. Из-за частичного сходства в постановке задач с классификацией кластеризацию также называют \textit{unsupervised classification}.

\subsection*{Классификация методов кластеризации}

По способу формирования кластеров выделяют следующие группы методов:

\begin{itemize}
    \item \textbf{Разделительные (partitioning)} — разбивают данные на заданное число кластеров, минимизируя расстояние внутри кластера и максимизируя расстояние между кластерами (например, K-means).
    
    \item \textbf{Основанные на плотности (density-based)} — группируют точки, находящиеся в областях с высокой плотностью, и отделяют их от областей с низкой плотностью (например, DBSCAN).
    
    \item \textbf{Основанные на сетке (grid-based)} — разбивают пространство на ячейки сетки и анализируют плотность данных в каждой ячейке (например, STING).
    
    \item \textbf{Основанные на модели (model-based)} — предполагают, что данные порождены некоторой статистической моделью и подбирают параметры этой модели (например, смеси Гауссианов).
    
    \item \textbf{Основанные на графах (graph-based)} — используют графовое представление данных и разбивают его на подграфы, соответствующие кластерам (например, спектральная кластеризация).
    
    \item \textbf{Основанные на подпространствах (subspace-based)} — ищут кластеры в подпространствах признаков, а не во всём пространстве (например, CLIQUE).
    
    \item \textbf{Основанные на ансамбле (ensemble-based)} — комбинируют результаты различных алгоритмов кластеризации для получения более стабильного и надёжного разбиения (например, CSPA).
\end{itemize}

\section*{Метод K-средних}

Кластеризация K-средних (\textit{K-means}) представляет собой одну из самых простых и популярных реализаций. Суть метода заключается в итеративной процедуре:

\begin{enumerate}
    \item Инициализация центроидов кластеров.
    \item Назначение каждой точки ближайшему центроиду.
    \item Пересчёт центроидов как среднего арифметического точек, принадлежащих кластеру.
    \item Повторение шагов 2--3 до сходимости.
\end{enumerate}

Оптимизация выполняется за счёт минимизации функции суммарного квадратичного отклонения точек от центроидов:

\[
J = \sum_{k=1}^{K} \sum_{x_i \in C_k} \|x_i - \mu_k\|^2
\]

где:
\begin{itemize}
    \item $K$ — число кластеров,
    \item $C_k$ — множество точек кластера $k$,
    \item $\mu_k$ — центр (центроид) кластера $k$.
\end{itemize}

\begin{figure}[H]
    \centering
    \includegraphics{clust.png}
    \caption{K-means}
\end{figure}

\section*{Описание датасета Seeds}

В работе используется датасет \textit{Seeds}, содержащий геометрические характеристики зёрен пшеницы. Каждый объект представляет собой отдельное зерно, описанное набором числовых признаков.

Датасет включает следующие признаки:

\begin{itemize}
    \item \textbf{area} — площадь зерна;
    \item \textbf{perimeter} — периметр зерна;
    \item \textbf{compactness} — компактность, рассчитываемая как отношение площади к квадрату периметра;
    \item \textbf{length\_kernel} — длина ядра зерна;
    \item \textbf{width\_kernel} — ширина ядра зерна;
    \item \textbf{asymmetry} — коэффициент асимметрии зерна;
    \item \textbf{groove\_length} — длина бороздки зерна;
    \item \textbf{class} — класс зерна (тип пшеницы).
\end{itemize}

Все признаки, кроме \textbf{class}, являются вещественными числовыми величинами и используются для кластеризации. Признак \textbf{class} представляет собой категориальную переменную, принимающую значения $1$, $2$ или $3$, и может использоваться для оценки качества кластеризации.


\endinput